\begin{frame}{Cellular Automata for Chloride Ingress}

\justifying

\textbf{Step 1: Aggregate Dataset Generation}

A DDPM is employed to generate a synthetic
aggregate dataset, capturing variability in size, shape, and spatial distribution
consistent with meso-scale representations of concrete.

\vspace{0.15cm}

\textbf{Step 2: Patch Construction}

The generated aggregates are assembled into spatial patches, which serve as the
computational domain for subsequent modeling of transport phenomena and material
deterioration.

\vspace{0.15cm}

\textbf{Step 3: Model Calibration}

The model is calibrated using the DE algorithm,
based on experimental chloride ingress data reported by \cite{DU201487}.

\vspace{0.15cm}

\textbf{Step 4: Numerical Simulation}

Simulations are performed on a concrete prism with dimensions
$75 \times 75$~mm to evaluate the chloride penetration front,
considering a concrete cover depth of 30~mm.

\end{frame}
